% GENERATED FILE. Edit canonical lessons, then run npm run generate:books.
% canonical-source-hash: fnv1a64:792ad6588ec02928
% canonical-lessons: HI-C548-avadhi, HI-C548-miyad, HI-C548-ant, HI-C548-alas, HI-C548-lutana

\chapter{Time Running Out}
\label{ch:hi-time-running-out}
\hlchaptermodality{548}

\begin{chapteropening}
\textbf{By the end of this chapter:} \emph{I can talk about time limits, endings and wasted chances.}

Complete the last lesson of chapter 548: I can talk about time limits, endings and wasted chances.
\end{chapteropening}

\section[\texorpdfstring{\hi{अवधि} (avadhi)}{अवधि (avadhi)}]{\hi{अवधि} (avadhi) — a period; a time limit}

\label{lesson:HI-C548-avadhi}



\begin{quote}
Before the new one: say the Hindi for an inkling, then the Hindi for an instruction.
\end{quote}

\subsection*{You'll want to know: \hi{अवधि}}
\textbf{\hi{अवधि}} — \emph{avadhi} — \textquotedblleft{}a period; a time limit\textquotedblright{}.

\begin{quote}
\hi{टिकट} \hi{की} \hi{अवधि} \hi{एक} \hi{हफ़्ता} \hi{है।}
\end{quote}

\begin{quote}
\emph{Ṭikaṭ kī avadhi ek haftā hai.}
\end{quote}

\begin{quote}
The ticket is valid for a week.
\end{quote}

Feminine.

\begin{cousinweb}
\textbf{\hi{अवधि}} (\emph{avadhi}) is a tatsama, taken whole from Sanskrit \emph{avadhi}, \textquotedblleft{}a term, a limit\textquotedblright{} (Monier-Williams).
\end{cousinweb}

\begin{grammarlens}[title={What you've built}]
One more word for saying what happens if.
\end{grammarlens}

\subsection*{Your turn}
\begin{itemize}
\raggedright
  \item \emph{Say it:} \emph{avadhi}
  \item \emph{Say it:} \emph{avadhi}, once more
  \item \emph{Say it:} \emph{Ṭikaṭ kī avadhi ek haftā hai.} — The ticket is valid for a week.
  \item \emph{Recall:} say the Hindi for an inkling, then the Hindi for an instruction, then say \emph{avadhi} again
\end{itemize}

\begin{tcolorbox}[breakable,colback=teal!4,colframe=teal!35!black,title={Before you move on}]
What does \hi{अवधि} mean? (\textquotedblleft{}A period; a time limit\textquotedblright{}.) Say it once more.
\end{tcolorbox}
\section[\texorpdfstring{\hi{मियाद} (miyād)}{मियाद (miyād)}]{\hi{मियाद} (miyād) — a term; an expiry date}

\label{lesson:HI-C548-miyad}



\begin{quote}
Before the new one: say the Hindi for an instruction, then the Hindi for a period.
\end{quote}

\subsection*{You'll want to know: \hi{मियाद}}
\textbf{\hi{मियाद}} — \emph{miyād} — \textquotedblleft{}a term; an expiry date\textquotedblright{}.

\begin{quote}
\hi{दवा} \hi{की} \hi{मियाद} \hi{दो} \hi{साल} \hi{है।}
\end{quote}

\begin{quote}
\emph{Davā kī miyād do sāl hai.}
\end{quote}

\begin{quote}
The medicine keeps for two years.
\end{quote}

Feminine.

\begin{cousinweb}
\textbf{\hi{मियाद}} (\emph{miyād}) is usually traced to Arabic \emph{mīʿād}, \textquotedblleft{}an appointed time\textquotedblright{}, through Persian.
\end{cousinweb}

\begin{grammarlens}[title={What you've built}]
One more word for saying what happens if.
\end{grammarlens}

\subsection*{Your turn}
\begin{itemize}
\raggedright
  \item \emph{Say it:} \emph{miyād}
  \item \emph{Say it:} \emph{miyād}, once more
  \item \emph{Say it:} \emph{Davā kī miyād do sāl hai.} — The medicine keeps for two years.
  \item \emph{Recall:} say the Hindi for an instruction, then the Hindi for a period, then say \emph{miyād} again
\end{itemize}

\begin{tcolorbox}[breakable,colback=teal!4,colframe=teal!35!black,title={Before you move on}]
What does \hi{मियाद} mean? (\textquotedblleft{}A term; an expiry date\textquotedblright{}.) Say it once more.
\end{tcolorbox}
\section[\texorpdfstring{\hi{अंत} (ant)}{अंत (ant)}]{\hi{अंत} (ant) — the end; an ending}

\label{lesson:HI-C548-ant}



\begin{quote}
Before the new one: say the Hindi for a period, then the Hindi for a term.
\end{quote}

\subsection*{You'll want to know: \hi{अंत}}
\textbf{\hi{अंत}} — \emph{ant} — \textquotedblleft{}the end; an ending\textquotedblright{}.

\begin{quote}
\hi{कहानी} \hi{का} \hi{अंत} \hi{बहुत} \hi{अच्छा} \hi{था।}
\end{quote}

\begin{quote}
\emph{Kahānī kā ant bahut acchā thā.}
\end{quote}

\begin{quote}
The story had a very good ending.
\end{quote}

Masculine.

\begin{cousinweb}
\textbf{\hi{अंत}} (\emph{ant}) is a tatsama, taken whole from Sanskrit \emph{anta}, \textquotedblleft{}end, limit\textquotedblright{} (Monier-Williams).
\end{cousinweb}

\begin{grammarlens}[title={What you've built}]
One more word for saying what happens if.
\end{grammarlens}

\subsection*{Your turn}
\begin{itemize}
\raggedright
  \item \emph{Say it:} \emph{ant}
  \item \emph{Say it:} \emph{ant}, once more
  \item \emph{Say it:} \emph{Kahānī kā ant bahut acchā thā.} — The story had a very good ending.
  \item \emph{Recall:} say the Hindi for a period, then the Hindi for a term, then say \emph{ant} again
\end{itemize}

\begin{tcolorbox}[breakable,colback=teal!4,colframe=teal!35!black,title={Before you move on}]
What does \hi{अंत} mean? (\textquotedblleft{}The end; an ending\textquotedblright{}.) Say it once more.
\end{tcolorbox}
\section[\texorpdfstring{\hi{आलस} (ālas)}{आलस (ālas)}]{\hi{आलस} (ālas) — laziness; sloth}

\label{lesson:HI-C548-alas}



\begin{quote}
Before the new one: say the Hindi for a term, then the Hindi for the end.
\end{quote}

\subsection*{You'll want to know: \hi{आलस}}
\textbf{\hi{आलस}} — \emph{ālas} — \textquotedblleft{}laziness; sloth\textquotedblright{}.

\begin{quote}
\hi{आलस} \hi{की} \hi{वजह} \hi{से} \hi{मेरी} \hi{बस} \hi{छूट} \hi{गई।}
\end{quote}

\begin{quote}
\emph{Ālas kī vajah se merī bas chūṭ gaī.}
\end{quote}

\begin{quote}
Out of laziness I missed my bus.
\end{quote}

Masculine.

\begin{cousinweb}
\textbf{\hi{आलस}} (\emph{ālas}) is a tadbhava, inherited from Sanskrit \emph{ālasya}, \textquotedblleft{}sloth\textquotedblright{} (CDIAL 1371).
\end{cousinweb}

\begin{grammarlens}[title={What you've built}]
One more word for saying what happens if.
\end{grammarlens}

\subsection*{Your turn}
\begin{itemize}
\raggedright
  \item \emph{Say it:} \emph{ālas}
  \item \emph{Say it:} \emph{ālas}, once more
  \item \emph{Say it:} \emph{Ālas kī vajah se merī bas chūṭ gaī.} — Out of laziness I missed my bus.
  \item \emph{Recall:} say the Hindi for a term, then the Hindi for the end, then say \emph{ālas} again
\end{itemize}

\begin{tcolorbox}[breakable,colback=teal!4,colframe=teal!35!black,title={Before you move on}]
What does \hi{आलस} mean? (\textquotedblleft{}Laziness; sloth\textquotedblright{}.) Say it once more.
\end{tcolorbox}
\section[\texorpdfstring{\hi{लुटाना} (luṭānā)}{लुटाना (luṭānā)}]{\hi{लुटाना} (luṭānā) — to squander; to give away freely}

\label{lesson:HI-C548-lutana}



\begin{quote}
Before the new one: say the Hindi for the end, then the Hindi for laziness.
\end{quote}

\subsection*{You'll want to know: \hi{लुटाना}}
\textbf{\hi{लुटाना}} — \emph{luṭānā} — \textquotedblleft{}to squander; to give away freely\textquotedblright{}.

\begin{quote}
\hi{उसने} \hi{अपने} \hi{पैसे} \hi{जुए} \hi{में} \hi{लुटा} \hi{दिए।}
\end{quote}

\begin{quote}
\emph{Usne apne paise jue meṁ luṭā diye.}
\end{quote}

\begin{quote}
He squandered his money on gambling.
\end{quote}

Past, with \hi{ने} (he or she …): \hi{उसने} \hi{पैसे} \hi{लुटाए}, \emph{usne paise luṭāe}.

\begin{cousinweb}
\textbf{\hi{लुटाना}} (\emph{luṭānā}) is the causative of \emph{lūṭnā}, \textquotedblleft{}to plunder\textquotedblright{}, which CDIAL derives from \emph{*luṭṭati}, a reconstructed form (CDIAL 11078): to let one's goods be carried off.
\end{cousinweb}

\begin{grammarlens}[title={What you've built}]
One more word for saying what happens if.
\end{grammarlens}

\subsection*{Your turn}
\begin{itemize}
\raggedright
  \item \emph{Say it:} \emph{luṭānā}
  \item \emph{Say it:} \emph{luṭānā}, once more
  \item \emph{Say it:} \emph{Usne apne paise jue meṁ luṭā diye.} — He squandered his money on gambling.
  \item \emph{Recall:} say the Hindi for the end, then the Hindi for laziness, then say \emph{luṭānā} again
\end{itemize}

\begin{tcolorbox}[breakable,colback=teal!4,colframe=teal!35!black,title={Before you move on}]
What does \hi{लुटाना} mean? (\textquotedblleft{}To squander; to give away freely\textquotedblright{}.) Say it once more.
\end{tcolorbox}
